\providecommand{\notesdir}{.}

% -----------------------------------------------------------------------------
% Shared preamble for the CO5059 topic notes, the final exam review and the
% cheat sheet's sibling documents.
% Documents outside this folder define \notesdir (path to this folder) first.
% -----------------------------------------------------------------------------

% Page layout and typography
\usepackage[a4paper,portrait,margin=1in]{geometry}
\usepackage{fontspec}
\usepackage{microtype}
\usepackage{parskip}
\setlength{\emergencystretch}{1.5em}% fewer overfull lines

% Mathematics and diagrams
\usepackage{amsmath,amssymb,amsthm,mathtools}
\usepackage{tikz-cd}
\usetikzlibrary{arrows.meta,positioning,shapes.geometric,shapes.arrows,
  shapes.multipart,calc,fit,backgrounds,matrix,patterns,decorations.pathreplacing}

% Reusable styles for flowcharts and concept diagrams.
\tikzset{
  concept/.style={
    draw=brandprimary,
    fill=brandsecondary!8,
    rounded corners,
    minimum width=3.2cm,
    minimum height=0.9cm,
    align=center,
    font=\small,
  },
  diagram arrow/.style={-{Latex[length=2.5mm]},thick,brandprimary},
  % tree nodes (e.g. Tree-LRU) and hardware blocks
  tnode/.style={draw,rectangle,font=\small\ttfamily,inner sep=3pt,
    minimum height=0.55cm,fill=white},
  box/.style={draw,rectangle,align=center,font=\small,inner sep=4pt,
    fill=white},
  % ALU drawn vertically (inputs on the top edge, output at the bottom):
  % \pic at (x,y) {valu}; inputs at (x-0.65,y), (x+0.65,y); output (x,y-1)
  valu/.pic={\draw[fill=white] (-1.1,0)--(-0.2,0)--(0,-0.3)--(0.2,0)--(1.1,0)
    --(0.5,-1)--(-0.5,-1)--cycle; \node[font=\scriptsize] at (0,-0.62) {ALU};},
}

% Tables, figures, code, and lists
\usepackage{graphicx}
\graphicspath{{\notesdir/}}
\usepackage{enumitem}
\usepackage{booktabs}
% array is pinned: a newer array.sty in the local texmf tree needs a newer
% LaTeX kernel than the installed one; the dated rollback works with both.
\usepackage{array}[=2024-06-01]
% P{w}: paragraph column set ragged right (no stretched word gaps in tables)
\newcolumntype{P}[1]{>{\raggedright\arraybackslash}p{#1}}
\usepackage{listings}

% Color and links
\usepackage[dvipsnames]{xcolor}
% Same palette as the other course notes (co5115):
% brandprimary = navy lines/arrows/links, brandsecondary!NN = light blue fills.
\definecolor{brandprimary}{HTML}{030391}
\definecolor{brandsecondary}{HTML}{1488D8}
\usepackage[
  colorlinks=true,
  linkcolor=brandprimary,
  citecolor=brandsecondary,
  urlcolor=brandsecondary
]{hyperref}

% Headers and footers
\usepackage{fancyhdr}
\setlength{\headheight}{26pt}
\pagestyle{fancy}
\fancyhf{}
\fancyhead[L]{\small\textbf{\coursecode}\\\coursename}
\fancyhead[R]{\small HCMUT}
\fancyfoot[L]{\small Semester \semester}
\fancyfoot[R]{\small Page \thepage}

% -----------------------------------------------------------------------------
% Course information -- change these values once per course
% -----------------------------------------------------------------------------
\newcommand{\universityname}{HCMC University of Technology}
\newcommand{\faculty}{Faculty of Computer Science and Engineering}
\newcommand{\coursecode}{CO5059}
\newcommand{\coursename}{Advanced System Architectures}
\newcommand{\semester}{261}
\newcommand{\instructor}{Assoc.\ Prof.\ Tran Ngoc Thinh}
\newcommand{\instructoremail}{}% leave empty to hide
\newcommand{\notetaker}{Phan Nhat Minh}
\newcommand{\notetakeremail}{minh.phannhat95@hcmut.edu.vn}
\newcommand{\emaillink}[1]{%
  \href{mailto:#1}{\textcolor{brandprimary}{\texttt{#1}}}}

% Definitions, results, examples, proofs and notes in coloured boxes.
\input{\notesdir/colored-boxes}

% -----------------------------------------------------------------------------
% Course-specific notation
% -----------------------------------------------------------------------------
\newcommand{\CPI}{\mathrm{CPI}}
\newcommand{\AMAT}{\mathrm{AMAT}}

\lstdefinestyle{plain}{
  basicstyle=\ttfamily\small,
  showstringspaces=false,
  columns=fullflexible,
  keepspaces=true,
  breaklines=true,
  frame=single,
  framerule=0.3pt,
  rulecolor=\color{black!30},
  xleftmargin=4pt,
  xrightmargin=4pt,
}
\lstdefinestyle{mips}{
  style=plain,
  keywordstyle=\bfseries,
  morekeywords={lw,sw,add,sub,and,or,slt,beq,bne,beqz,xor,nop,addi,addiu,addu,
    subu,li,la,jal,jr,jalr,sll,srl,sra,sllv,srlv,srav,andi,ori,xori,nor,slti,
    sltiu,sltu,lui,mult,multu,div,divu,mfhi,mflo,ll,sc,syscall,move,lh,lhu,lb,
    lbu,sh,sb,blez,bgtz,bltz,bgez,j,mul,ld,sd,dadd,daddu,daddi,daddui,dsub,
    dsubu,bnez,l.d,s.d,add.d,sub.d,mul.d,div.d},
}
\lstset{style=mips}

% Long identifiers (GRADE_REPORT.Section_identifier) may break after "_".
\makeatletter
\DeclareRobustCommand{\_}{\ifmmode\nfss@text{\textunderscore}\else
  \textunderscore\allowbreak\fi}
\makeatother

% Marks a correction of an error in the slides: \begin{remark}\slidefix{} ...
\newcommand{\slidefix}{\textbf{\textcolor{red!75!black}{Slide correction.}}}
% Marks material added from the textbook or papers, not in the slides.
\newcommand{\beyond}{\textbf{Beyond the slides.}}

% -----------------------------------------------------------------------------
% Textbook citations as footnotes (book page + PDF page of textbooks/*.pdf)
% -----------------------------------------------------------------------------
% \srcnote{text}: a footnote that also works inside the coloured boxes
% (mdframed/tcolorbox would otherwise keep it inside the box with a letter).
% In a box it prints the mark and queues the text; the queue is emptied after
% the box (hooks below), so the note lands at the foot of the page.
\makeatletter
\newtoks\srcnote@pending
\def\srcnote@mp{mpfootnote}
\newcommand{\srcnote}[1]{%
  \ifx\@mpfn\srcnote@mp
    \stepcounter{footnote}%
    \protected@xdef\@thefnmark{\thefootnote}%
    \@footnotemark
    \xdef\srcnote@n{\the\c@footnote}%
    \global\srcnote@pending\expandafter{\the\expandafter\srcnote@pending
      \expandafter\footnotetext\expandafter[\srcnote@n]{#1}}%
  \else\footnote{#1}\fi}
\newcommand{\srcnote@flush}{\the\srcnote@pending\global\srcnote@pending{}}
\@for\srcnote@env:={theorem,definition,proposition,prop,lemma,corollary,%
exercise,tbexercise,example,eg,remark,note,solution,explanation,noot}\do{%
  \AddToHook{env/\srcnote@env/after}{\srcnote@flush}}
\makeatother
% Hennessy & Patterson, Computer Architecture: A Quantitative Approach, 5th ed.
% (references/Computer Architecture - A Quantitative Approach.pdf).
% Chapters: PDF page = printed page + 29.  Appendix X: page X-n is PDF n + offset.
\makeatletter
\@namedef{HPoff@}{29}
\@namedef{HPoff@A}{523}\@namedef{HPoff@B}{579}\@namedef{HPoff@C}{647}
\@namedef{HPoff@D}{736}\@namedef{HPoff@E}{804}\@namedef{HPoff@F}{831}
\@namedef{HPoff@G}{949}\@namedef{HPoff@H}{984}\@namedef{HPoff@I}{1029}
\@namedef{HPoff@J}{1076}\@namedef{HPoff@K}{1150}
% \hppage{X}{n}: "p. X-n (PDF p. m)"; X empty for a chapter page.
\newcommand{\hppage}[2]{p.~\if\relax\detokenize{#1}\relax\else#1-\fi#2\
  (PDF~\the\numexpr#2+\csname HPoff@#1\endcsname\relax)}
% \hppages{X}{first}{last}
\newcommand{\hppages}[3]{pp.~\if\relax\detokenize{#1}\relax\else#1-\fi#2--%
  \if\relax\detokenize{#1}\relax\else#1-\fi#3\
  (PDF~\the\numexpr#2+\csname HPoff@#1\endcsname\relax--%
  \the\numexpr#3+\csname HPoff@#1\endcsname\relax)}
\makeatother
% Footnote citations:
%   \HP{location}{page}            chapter page          \HP[last]{loc}{first}
%   \HPA{X}{location}{page}        appendix X page X-n   \HPA[last]{X}{loc}{first}
\newcommand{\HP}[3][]{\srcnote{H\&P 5th ed., #2,
  \if\relax\detokenize{#1}\relax\hppage{}{#3}\else\hppages{}{#3}{#1}\fi.}}
\newcommand{\HPA}[4][]{\srcnote{H\&P 5th ed., #3,
  \if\relax\detokenize{#1}\relax\hppage{#2}{#4}\else\hppages{#2}{#4}{#1}\fi.}}

% Textbook exercise. No solution is written: the box cites the exercise and
% the book pages with the theory needed to solve it.
% \begin{tbexercise}[appendix letter]{H&P number}{page}{theory pointer}
%   e.g. \begin{tbexercise}{1.8}{60}{\S1.9, \hppages{}{46}{48}}
%        \begin{tbexercise}[B]{B.1}{60}{\S B.2, \hppages{B}{15}{22}}
\makeatletter
\newenvironment{tbexercise}[4][]
  {\gdef\tbx@note{\srcnote{H\&P 5th ed., Exercise #2, \hppage{#1}{#3}.}}%
   \gdef\tbx@theory{#4}%
   \begin{exercise}[H\&P #2]}
  {\par\hfill{\small\itshape Theory: \tbx@theory}\tbx@note
   \end{exercise}}
\makeatother
% Source note under a diagram, for checking it against the original:
% \diagramsource{slides file, slide N; cf. textbook figure}
\newcounter{diagram}
% (the counter is stepped inside the note's line, so no page break can fall
% between the diagram and its note)
\newcommand{\diagramsource}[1]{\par\nopagebreak
  {\centering\footnotesize\itshape\color{black!65}\leavevmode\refstepcounter{diagram}%
   Diagram~\thediagram\ -- source: #1\par}}

% Source line under a section heading.
\newcommand{\sources}[1]{\par\noindent{\small\itshape\sloppy Sources: #1\par}\medskip}

% Cover page + contents page.  Usage: \makecover{Big title}
\newcommand{\makecover}[1]{%
\pagenumbering{roman}
\begin{titlepage}
  \thispagestyle{empty}
  \renewcommand{\emaillink}[1]{\href{mailto:##1}{\textcolor{black}{\texttt{##1}}}}% no colour on the cover

  \begin{center}
    \includegraphics[width=0.62\textwidth]{assets/text_logo_3.png}

    \vspace{0.35cm}
    {\large\bfseries\faculty\par}
  \end{center}

  \vfill
  \begin{center}
    {\Huge\bfseries #1\par}
    \vspace{0.55cm}
    {\LARGE\bfseries\coursename\par}
    \vspace{0.2cm}
    {\large\coursecode\par}

    \vspace{0.65cm}
    \rule{0.38\textwidth}{0.6pt}
    \vspace{0.65cm}

    {\large\textbf{Instructor:} \instructor
      \ifx\instructoremail\empty\else\ {\small(\emaillink{\instructoremail})}\fi\par}
    \vspace{0.3cm}
    {\large\textbf{Notes by:} \notetaker\
      {\small(\emaillink{\notetakeremail})}\par}
    \vspace{0.3cm}
    {\large\textbf{Semester:} \semester\par}
  \end{center}

  \vfill
  \begin{center}
    \small Last updated: \today
  \end{center}
\end{titlepage}

% -----------------------------------------------------------------------------
% Table of contents -- kept on its own page
% -----------------------------------------------------------------------------
\clearpage
\setcounter{page}{2}
\thispagestyle{empty}
\hypersetup{linkcolor=black}
\setcounter{tocdepth}{2}
\tableofcontents
\hypersetup{linkcolor=brandprimary}
\clearpage
\pagenumbering{arabic}
}
